\documentclass[a4paper,11pt]{article}

\usepackage{fontspec}
\setmainfont{XCharter}[
  UprightFont = *-Roman,
  BoldFont = *-Bold,
  ItalicFont = *-Italic,
  BoldItalicFont = *-BoldItalic,
  Ligatures = TeX,
]
\usepackage[a4paper,left=30pt,right=30pt,top=24pt,bottom=24pt]{geometry}
\usepackage{titlesec}
\usepackage{enumitem}
\usepackage[hidelinks]{hyperref}
\usepackage{fancyhdr}
\usepackage{xcolor}
\usepackage{multicol}

\definecolor{accent}{HTML}{1A2733}
\definecolor{subtle}{HTML}{666666}

\pagestyle{fancy}
\fancyhf{}
\renewcommand{\headrulewidth}{0pt}
\renewcommand{\footrulewidth}{0pt}

\urlstyle{rm}
\raggedright
\linespread{0.9}
\setlength{\parindent}{0pt}
\setlength{\tabcolsep}{0in}

\titleformat{\section}{%
  \normalfont\large\bfseries\color{accent}%
}{}{0em}{}[{\color{subtle}\vspace{1pt}\titlerule[0.4pt]}]
\titlespacing*{\section}{0pt}{6pt}{2pt}

\newcommand{\resumeSubheading}[4]{%
  \vspace{1pt}%
  \noindent{\bfseries #1}\hfill {\small #4}\par%
  \noindent{\itshape\color{subtle} #3 \,$\cdot$\, #2}\par%
}

\newcommand{\resumeProject}[3]{%
  \vspace{1pt}%
  \if\relax\detokenize{#3}\relax
  \noindent{\bfseries #1}\quad{\itshape\color{subtle} #2}\par%
  \else
  \noindent{\bfseries #1}\hfill {\small #3}\par%
  \noindent{\itshape\color{subtle} #2}\par%
  \fi
}

\setlist[description]{leftmargin=14pt, itemsep=1pt, topsep=1pt, parsep=0pt}
\setlist[itemize]{leftmargin=14pt, label=\textbullet, itemsep=0pt, topsep=0pt, parsep=0pt}

\begin{document}

%----------HEADING-----------------
\begin{center}
  {\fontsize{24pt}{28pt}\selectfont\bfseries\color{accent} Mustafa Arinmis}\\[4pt]
  {\small\color{subtle} mustafa\_arinmis@outlook.com \, $\bullet$ \, +90 543 549 99 53 \, $\bullet$ \, \href{https://github.com/arinmis}{github.com/arinmis} \, $\bullet$ \, \href{https://www.linkedin.com/in/arinmis}{linkedin.com/in/arinmis}}\\[2pt]
  {\small\color{subtle} Antalya, Türkiye \,$\cdot$\, Bilgisayar Mühendisliği, Akdeniz Üniversitesi}
\end{center}
\vspace{2pt}

%-----------SUMMARY-----------------
\section{Özet}
\textbf{4 yıllık yazılım mühendisliği deneyimine} sahip bir \textbf{Yapay Zeka Mühendisiyim}. Backend, mikroservis ve full-stack geliştirme geçmişimi AI engineering yetkinlikleriyle birleştiriyorum. \textbf{Agentic sistemler}, \textbf{RAG pipeline'ları} ve \textbf{agent harness} geliştiriyorum. Agent orchestration, tool calling ve MCP ile LLM uygulamalarını production sistemlere entegre ediyorum.

\section{Deneyim}
\resumeSubheading{\href{https://www.createanapi.com}{createanapi.com} (\href{https://www.promolink.io/}{PromoLink})}{Uzaktan, Arizona, ABD}{Full Stack Yapay Zeka Yazılım Mühendisi}{Ağustos 2026 -- Devam ediyor}
\begin{itemize}
  \item Promosyon ürünlerine yönelik \textbf{e-ticaret platformu PromoLink'te} .NET/C#, React, AWS ve MongoDB ile full-stack özellikler ve AI agent'lar geliştiriyorum.
  \item \textbf{400.000'den fazla ürün içeren katalogda agentic RAG tabanlı arama} geliştiriyorum; OpenSearch semantic retrieval ve keyword search ile fiyat, adet ve malzeme filtrelerini birleştiriyorum.
  \item Doğal dildeki taleplerle online mağaza oluşturan ve düzenleyen \textbf{CMS agent} geliştiriyorum; tool calling ve schema validation kullanıyor, canlı preview sunuyorum.
  \item Bu özellikleri destekleyen \textbf{LLM'den bağımsız agent platformunu} geliştiriyorum; .NET tool handler'ları, retry policy, event recovery, embedding batching/caching ve concurrent retrieval kullanıyorum.
\end{itemize}

\resumeSubheading{\href{https://www.santsg.com/en/products/octopus-project/}{SAN TSG (OctopusPlus)}}{Antalya, Türkiye}{Backend Yazılım Mühendisi}{Kasım 2024 -- Temmuz 2026}
\begin{itemize}
  \item 30'dan fazla .NET mikroservisi için runtime tool discovery destekleyen JWT korumalı MCP server ve LLM agent platformu geliştirdim.
  \item Teknik dokümantasyon için \textbf{RAG servisi} geliştirdim (Flask, ChromaDB); 6 stajyere mentorluk yaptım, biri yarı zamanlı mühendis olarak işe alındı.
  \item Turizm rezervasyonları için \textbf{5 production mikroservisi} .NET/C#, CQRS, Dapr ve RabbitMQ ile geliştirip AWS EKS üzerinde canlıya aldım; GitOps için ArgoCD kullandım.
\end{itemize}

\resumeSubheading{P.I. Works}{Uzaktan, İstanbul, Türkiye}{Full Stack Yazılım Geliştirici}{Haziran 2023 -- Eylül 2024}
\begin{itemize}
  \item Kurumsal coğrafi konum platformunda Angular ve C#/.NET ile özellikler geliştirdim, birim test kapsamını artırdım ve kodun bakımını kolaylaştırdım.
  \item Ağ sorunlarına yönelik müdahaleleri izleyen kök neden analizi araçları geliştirdim; \textbf{Docker} ile paketleyip \textbf{GitLab CI} üzerinden dağıttım.
\end{itemize}

\resumeSubheading{Linarit}{Antalya, Türkiye}{Frontend Geliştirici}{Aralık 2021 -- Eylül 2022}
\begin{itemize}
  \item TÜBİTAK Ar-Ge projesinde endüstriyel IoT dashboard'unun React frontend'ini üstlendim; gerçek zamanlı gRPC telemetri görselleştirmeleri, rendering optimizasyonu ve client-side caching geliştirdim.
\end{itemize}

\section{Eğitim}
\resumeProject{Bilgisayar Mühendisliği, Yüksek Lisans}{Akdeniz Üniversitesi, Antalya}{2026 -- Devam ediyor}
\resumeProject{Bilgisayar Mühendisliği, Lisans}{Akdeniz Üniversitesi, Antalya}{2019 -- 2024}

\section{Teknik Yetkinlikler}
\textbf{AI engineering:} Agent orchestration, tool calling, RAG, embeddings, MCP, semantic search\\
\textbf{Diller ve framework'ler:} C\#, .NET, Python, React, TypeScript, Flask, FastAPI\\
\textbf{Veri ve altyapı:} OpenSearch, MongoDB, ChromaDB, PostgreSQL, Redis, Docker, AWS EKS/ECR

\section{Projeler}
\resumeProject{Tiro}{\href{https://tiro.legal}{tiro.legal}}{}
\begin{itemize}
  \item Avukatlar için cross-platform, RAG tabanlı dava yönetim uygulaması geliştirdim; \textbf{local LLM runtime} entegre edip model fine-tuning yaptım. Python, FastAPI, ChromaDB, Electron, React.
\end{itemize}

\resumeProject{FitBrand}{\href{https://fitbrand.app}{fitbrand.app}}{}
\begin{itemize}
  \item AI destekli antrenman planları ve MCP tools içeren white-label antrenör platformu. Next.js, Prisma, PostgreSQL; markaya özel PWA ve danışan yönetimi.
\end{itemize}
\end{document}
